\documentclass[11pt]{article}
\usepackage[T1]{fontenc}
\usepackage{lmodern}
\usepackage[margin=1in]{geometry}
\usepackage{amsmath,amssymb,amsthm}
\usepackage{microtype}
\usepackage[hidelinks]{hyperref}
\newtheorem{theorem}{Theorem}
\newtheorem{lemma}[theorem]{Lemma}
\newcommand{\rev}{\mathcal R}
\title{A missing multiplicity condition for double-permutation sequences}
\author{Alec Kriebel\\
  \small Independent researcher\\
  \small ORCID: \href{https://orcid.org/0009-0001-9320-500X}{0009-0001-9320-500X}}
\date{16 October 2026\\ \small Research note, version 1.0}
\begin{document}
\maketitle
\begin{abstract}
The four term-sequence axioms in the inspected Goodman--Pollack manuscript
and the 2008 Oberwolfach problem session do not suffice for the proposed
encoding by connected sets with compatible tangent and separating
pseudolines. An explicit two-body word has least period eight and satisfies
all four axioms, but has no separating switch and hence cannot encode
the required internal tangents. This gives a negative answer to the literal
universal realization question under those axioms. It leaves intact the
representation theorem for valid fibrations and the stronger
allowable-interval definition. We also prove that requiring a fully
separated term forces the missing switch multiplicities, clarifying the
boundary of a separate normalized-input construction request.
\end{abstract}

\noindent\textbf{Draft status.}
Mathematical verification is complete; the independent priority audit and
whole-package publication reviews remain pending. This draft is not ready
for submission or deposit.

\medskip
\noindent\textbf{Disclosure.}
AI tools were used extensively in formulating, reproducing, checking and
adversarially reviewing the argument, investigating sources, and preparing
this note and its software. This is an unrefereed preprint. Internal
AI-assisted review is not external human peer review or proof-assistant
certification.

\section{The exact formulation}

Fix an integer $n\ge2$. A double permutation is an ordering of
\[
  1,1',\ldots,n,n'
\]
in which $i$ precedes $i'$ for every $i$. Let $\rev$ reverse an ordering
and interchange every primed and unprimed index. Set $N=8\binom n2$.
Consider a doubly infinite word $(P_k)_{k\in\mathbb Z}$ satisfying:
\begin{enumerate}
\item Each $P_k$ is a double permutation.
\item Each transition $P_k\longrightarrow P_{k+1}$ swaps two adjacent
indices belonging to different bodies.
\item $P_{k+N/2}=\rev(P_k)$ for every $k$.
\item The word has period $N$.
\end{enumerate}
We call these the \emph{four term axioms}. The example below has
\emph{least} period $N$, so it also satisfies the interpretation of the
last axiom as a primitive-period requirement.

These are Corollary~11 and Definition~12 of the 26 November 2006
Goodman--Pollack author manuscript \cite[pp.~8--9]{GP06}; the same four
conditions are printed in Pollack's problem-session item in
\cite[p.~2551]{OWR}. The topological realization question appears explicitly
in the manuscript and as Question~1 in \cite[p.~2492]{OWR}.
Our input definition is exactly the four conditions above.

The source geometric encoding uses pairwise disjoint compact connected
sets and a compatible arrangement of pseudolines, with a strict separator
and four designated common tangents for each body pair: two internal and
two external \cite[Theorem~1 and Definitions~4,~8--9]{GP06}.
Both directions of every tangent are recorded. A directed tangent is
indexed according to the left or right supporting side of each body.
Thus an internal tangent, with the bodies on opposite sides, gives an
ordered switch involving one primed and one unprimed index. We call
such a switch \emph{mixed}. Supporting switches have two primed or two
unprimed indices. We use the post-switch ordering convention of
\cite[p.~7]{GP06}.

\begin{theorem}\label{thm:negative}
The four term axioms do not imply realizability in this connected-set,
compatible-pseudoline encoding. A counterexample already exists for
$n=2$, with least period eight.
\end{theorem}

\section{The eight-term obstruction}

Put
\[
 A=(1,2,1',2'),\qquad B=(1,2,2',1'),\qquad C=(2,1,2',1').
\]
Repeat the following period in both directions:
\[
 \boxed{A,\ B,\ A,\ B,\ C,\ B,\ C,\ B.}                 \tag{1}
\]
Every displayed ordering has $1$ before $1'$ and $2$ before $2'$.
Each of the transitions $A\leftrightarrow B$ and
$B\leftrightarrow C$ is one adjacent different-body swap. These are all
the transitions of (1), including its wrap.

Moreover $\rev(A)=C$, $\rev(C)=A$ and $\rev(B)=B$, so the half-period
identity holds at every phase. The least period divides eight. A proper
divisor would be $1$, $2$ or $4$ and would imply $P_0=P_4$.
Since $A\ne C$, the least period is eight.

The ordered post-switch pairs, including the last-to-first transition,
are
\[
\begin{gathered}
 (2',1'),\ (1',2'),\ (2',1'),\ (2,1),\\
 (1,2),\ (2,1),\ (1,2),\ (1',2').
\end{gathered}                                                   \tag{2}
\]
Each of the four supporting types occurs twice. None of
$(1,2')$, $(2,1')$, $(1',2)$ or $(2',1)$ occurs.

In the source encoding the two internal tangent pseudolines of this
body pair must contribute mixed events. In fact recording both directions
gives four such events. Equation (2) has none, so it cannot be that
encoding. This proves Theorem~\ref{thm:negative}.

This obstruction uses the prescribed supporting sides and tangent roles,
rather than Euclidean projection or simultaneous straightening. The
source's strict separators and compatibility hypotheses ensure the
unambiguous directed-tangent encoding; the argument requires neither
smoothness nor positive area. A degeneracy that does not supply the
prescribed tangent roles and single-swap events does not realize (1).
Changing phase, reversing traversal, renaming the bodies or switching
between before- and after-swap conventions preserves the absence of
mixed events.

\section{A fully separated input term changes the problem}

The last paragraph of the separate polygonal-construction request
\cite[p.~2551]{OWR} refers to the fully separated ordering
$(1,1',2,2',\ldots,n,n')$. A variant that explicitly requires this
ordering in its input excludes (1). More generally:

\begin{lemma}\label{lem:separated}
If a word satisfying the four term axioms contains
\[
 D=(i_1,i_1',i_2,i_2',\ldots,i_n,i_n')
\]
for some ordering $i_1,\ldots,i_n$ of the body labels, then every
ordered different-body endpoint pair occurs exactly once as a switch
in a full period.
\end{lemma}

\begin{proof}
Shift the phase to make $P_0=D$. Its half-period counterpart
$\rev(D)$ reverses the order of the body blocks and preserves each
body's internal order. The relative order of each of the four endpoint
pairs belonging to any two different bodies must therefore reverse.
There are exactly $4\binom n2=N/2$ such pairs.

An adjacent transposition changes the relative order of exactly the
two endpoints it swaps. Each required endpoint pair must consequently
be swapped an odd, positive number of times between $D$ and $\rev(D)$.
There are precisely $N/2$ steps available, so every required pair is
swapped once and none can backtrack. Reverse-toggle transports the
second half to the first, giving one swap of each unordered
different-body endpoint pair in each half. Since the full period
returns every relative order to its starting value, the two swaps of
each such pair have opposite orders. Thus every ordered type appears
once in the full period.
\end{proof}

The word (1) has no fully separated term. Lemma~\ref{lem:separated}
is a multiplicity statement, not a construction of polygons or a
realization theorem in higher $n$. The four term axioms alone do not
require a fully separated term. The reference to that term in the
polygonal request must therefore be distinguished from an explicit
normalized-input hypothesis. We do not prove the additional polygonal
construction requirements in either interpretation.

\section{Literature and verification boundaries}

The stronger definition was already explicit in
Dhandapani--Goodman--Holmsen--Pollack \cite[Section~2, p.~5]{DGHP}:
each directed separating and supporting type occurs once per full
period. Novick's published definition \cite[pp.~1060--1061]{Novick}
requires two supporting and two separating switches per body pair
per half-period. Both definitions exclude (1). The missing
multiplicity condition is established literature, not a new definition.

The Oberwolfach contribution immediately attributes an affirmative
answer to Habert and Pocchiola \cite[pp.~2492--2493]{OWR}.
Their final published Theorem~49 \cite[p.~624]{HP} realizes an already
specified valid fibration of an arrangement of pseudocircles.
Its following application identifies the properly constrained interval
encodings with fibrations. Theorem~\ref{thm:negative} concerns the weaker
four term axioms; it supplies no counterexample to that representation
theorem. We do not independently reprove the representation theorem or
the full coding equivalence for the stronger input class.

The contribution here is the explicit finite obstruction to the literal
four-axiom universal assertion and the boundary lemma above. The inspected
definitions are those of the dated author manuscript and the official
2008 report. The final Goodman--Pollack journal body was not obtained;
no assertion about a correction or its absence in that edition is made.
The documented priority search is bounded. It certifies neither historical
firstness of this observation nor that a corresponding question remains
open today.

The accompanying source and verification archive preserves the frozen
candidate artifacts and independent reviews. Separate exact enumerations
check all $6^4=1296$ possible first halves for two bodies. They agree on
48 rooted primitive weak-axiom words: 16 satisfy switch-once and 32 have
no mixed event. These are indexed words, not isomorphism-class counts.
The archive also checks invalid inputs, the symmetry transforms of (1),
positive separated-start controls and artifact integrity. The
counterexample and Lemma~\ref{lem:separated} do not depend on the enumeration.
Finite software checks do not certify a topological representation proof,
literature completeness or formal proof-assistant verification.

\begin{thebibliography}{9}

\bibitem{GP06}
J.~E. Goodman and R. Pollack,
\emph{The combinatorial encoding of disjoint convex sets in the plane},
author manuscript, 26 November 2006, 13 pp.
\url{https://math.nyu.edu/inmemoriam/pollack/pubs/dblperm11-26-06.pdf}.
References here use this manuscript's numbering.

\bibitem{OWR}
R. Pollack,
\emph{Allowable double permutation sequences, arrangements of double
pseudolines of Michel Pocchiola and their applications to planar families
of convex sets}, pp.~2491--2493; and
\emph{Allowable double-permutation sequences}, p.~2551,
in \emph{Discrete Geometry}, Oberwolfach Report 44/2008.
\href{https://doi.org/10.4171/OWR/2008/44}{doi:10.4171/OWR/2008/44}.

\bibitem{DGHP}
R. Dhandapani, J.~E. Goodman, A. Holmsen and R. Pollack,
\emph{Interval sequences and the combinatorial encoding of planar
families of pairwise disjoint convex sets},
author manuscript, 21 April 2005, 19 pp., Section~2.
\url{https://math.nyu.edu/inmemoriam/pollack/pubs/interval.pdf}.
The journal citation is Rev. Roumaine Math. Pures Appl.
\textbf{50} (2005), 537--553; page references here use the manuscript.

\bibitem{Novick}
M. Novick,
\emph{Allowable Interval Sequences and Line Transversals in the Plane},
Discrete Comput. Geom. \textbf{48} (2012), 1058--1073.
\href{https://doi.org/10.1007/s00454-012-9457-x}{doi:10.1007/s00454-012-9457-x}.

\bibitem{HP}
L. Habert and M. Pocchiola,
\emph{LR Characterization of Chirotopes of Finite Planar Families of
Pairwise Disjoint Convex bodies},
Discrete Comput. Geom. \textbf{50} (2013), 552--648.
\href{https://doi.org/10.1007/s00454-013-9532-y}{doi:10.1007/s00454-013-9532-y}.
The fibration definition and Theorem~49 are on pp.~623--624.

\end{thebibliography}
\end{document}
